% !TEX root = z-main.tex

El proyecto tuvo como propósito preservar digitalmente y divulgar el patrimonio cultural guatemalteco mediante el desarrollo de modelos tridimensionales del sitio arqueológico Kaminaljuyú. Este sitio, de gran relevancia histórica para el altiplano central de Guatemala, se encuentra en gran parte sepultado por la urbanización. Esta situación ha dificultado su acceso y estudio.

Para atender esta problemática, se elaboraron seis modelos tridimensionales de montículos seleccionados a partir de información arqueológica validada por especialistas. El trabajo permitió reconstruir digitalmente estructuras que no son visibles en la actualidad y preparar dichos modelos para su integración en una aplicación de realidad aumentada orientada a la educación y la difusión cultural.

Se obtuvieron representaciones tridimensionales fieles. Estas incorporaron texturas y acabados que favorecieron la comprensión de la apariencia original de los montículos. Los archivos finales se almacenaron en formatos compatibles, garantizando su preservación y posible reutilización en futuros proyectos de investigación o divulgación.

Se concluyó que la digitalización de sitios arqueológicos mediante modelado tridimensional y su aplicación en entornos de realidad aumentada constituye una estrategia eficaz para la conservación del patrimonio. Se sugiere su continuidad en otras áreas arqueológicas del país.